\section{Experimental research}

This chapter describes the experiments conducted to evaluate the proposed system.

\subsection{Experimental setup}

The experiments were performed on the following hardware:
\begin{itemize}
  \item CPU: 8-core processor, 3.2~GHz;
  \item RAM: 32~GB DDR4;
  \item Storage: 512~GB NVMe SSD;
  \item OS: Ubuntu 24.04 LTS.
\end{itemize}

\subsection{Test datasets}

Three datasets of varying sizes were used for evaluation:

% === EXAMPLE: Table with multirow cells ===
\begin{table}[H]
    \caption{Test datasets} \label{tab:datasets}
    \begin{tabular}{|l|c|c|c|}
    \hline \textbf{Dataset} & \textbf{Records} & \textbf{Size (MB)} & \textbf{Description} \\
    \hline Small   & 1\,000    & 5   & Unit testing \\
    \hline Medium  & 100\,000  & 450 & Integration testing \\
    \hline Large   & 1\,000\,000 & 4\,200 & Stress testing \\
    \hline
    \end{tabular}
\end{table}

\subsection{Performance results}

The system was benchmarked against two baseline implementations.
Table~\ref{tab:performance} shows the results.

% === EXAMPLE: Table with multirow and merged cells ===
\begin{table}[H]
    \caption{Performance comparison (processing time in seconds)} \label{tab:performance}
    \begin{tabular}{|l|c|c|c|c|}
    \hline
    \multirow{2}{*}{\textbf{Method}} & \multicolumn{3}{c|}{\textbf{Dataset size}} & \multirow{2}{*}{\textbf{Avg. improvement}} \\
    \cline{2-4}
      & \textbf{Small} & \textbf{Medium} & \textbf{Large} & \\
    \hline Baseline A       & 0.8  & 45.2  & 512.0 & --- \\
    \hline Baseline B       & 0.5  & 38.7  & 478.3 & --- \\
    \hline Proposed approach & 0.3  & 22.1  & 287.5 & 40\% \\
    \hline
    \end{tabular}
\end{table}

The proposed approach demonstrates a \textbf{40--44\%} improvement in processing time compared to Baseline~A, and a \textbf{37--40\%} improvement over Baseline~B.

% === EXAMPLE: TikZ bar chart (manual) ===
Figure~\ref{fig:bar-chart} visualizes the performance comparison across dataset sizes.

\begin{figure}[H]
    \centering
    \begin{tikzpicture}[
        bar/.style={minimum width=0.5cm, anchor=south},
        >=Stealth
    ]
        % Axes
        \draw[->] (0,0) -- (9.5,0) node[right] {Dataset};
        \draw[->] (0,0) -- (0,5.5) node[above] {Time (s)};

        % Scale: 1cm = 100s
        \foreach \y/\val in {1/100, 2/200, 3/300, 4/400, 5/500} {
            \draw (-0.1,\y) -- (0.1,\y) node[left=2pt] {\val};
        }

        % Small dataset (scaled x10 for visibility)
        \node at (1.5, -0.4) {\small Small};
        \fill[red!60] (0.8,0) rectangle (1.2,0.08);
        \fill[blue!60] (1.3,0) rectangle (1.7,0.05);
        \fill[green!60] (1.8,0) rectangle (2.2,0.03);

        % Medium dataset
        \node at (4.5, -0.4) {\small Medium};
        \fill[red!60] (3.8,0) rectangle (4.2,0.452);
        \fill[blue!60] (4.3,0) rectangle (4.7,0.387);
        \fill[green!60] (4.8,0) rectangle (5.2,0.221);

        % Large dataset
        \node at (7.5, -0.4) {\small Large};
        \fill[red!60] (6.8,0) rectangle (7.2,5.12);
        \fill[blue!60] (7.3,0) rectangle (7.7,4.783);
        \fill[green!60] (7.8,0) rectangle (8.2,2.875);

        % Legend
        \fill[red!60] (9.5,4.5) rectangle (10,4.8);
        \node[right] at (10,4.65) {\small Baseline A};
        \fill[blue!60] (9.5,3.8) rectangle (10,4.1);
        \node[right] at (10,3.95) {\small Baseline B};
        \fill[green!60] (9.5,3.1) rectangle (10,3.4);
        \node[right] at (10,3.25) {\small Proposed};
    \end{tikzpicture}
    \caption{Performance comparison across dataset sizes} \label{fig:bar-chart}
\end{figure}

\subsection{Accuracy results}

Accuracy was measured using standard metrics:

\begin{equation} \label{eq:accuracy}
    \text{Accuracy} = \frac{TP + TN}{TP + TN + FP + FN} \times 100\%
\end{equation}

where $TP$, $TN$, $FP$, $FN$ denote true positives, true negatives, false positives, and false negatives respectively.

% === EXAMPLE: Additional classification metrics ===
The precision, recall, and F1-score are defined as:

\begin{align}
    \text{Precision} &= \frac{TP}{TP + FP} \label{eq:precision} \\[6pt]
    \text{Recall} &= \frac{TP}{TP + FN} \label{eq:recall} \\[6pt]
    F_1 &= 2 \cdot \frac{\text{Precision} \cdot \text{Recall}}{\text{Precision} + \text{Recall}} \label{eq:f1}
\end{align}

The proposed approach achieved an accuracy of \textbf{94.2\%}, compared to 89.5\% for Baseline~A and 91.1\% for Baseline~B.

% === EXAMPLE: Detailed metrics table with multirow ===
\begin{table}[H]
    \caption{Detailed classification metrics} \label{tab:metrics}
    \begin{tabular}{|l|c|c|c|c|}
    \hline
    \multirow{2}{*}{\textbf{Method}} & \multicolumn{4}{c|}{\textbf{Metrics (\%)}} \\
    \cline{2-5}
      & \textbf{Accuracy} & \textbf{Precision} & \textbf{Recall} & \textbf{$F_1$-score} \\
    \hline Baseline A       & 89.5 & 87.3 & 88.1 & 87.7 \\
    \hline Baseline B       & 91.1 & 90.2 & 89.8 & 90.0 \\
    \hline Proposed approach & \textbf{94.2} & \textbf{93.5} & \textbf{93.8} & \textbf{93.6} \\
    \hline
    \end{tabular}
\end{table}

\subsection{Scalability analysis}

The system scales linearly with dataset size, as evidenced by the following relationship:

\begin{equation} \label{eq:scalability}
    T(n) = a \cdot n + b
\end{equation}

where $T(n)$ is the processing time, $n$ is the number of records, and $a$, $b$ are constants determined through regression analysis ($a = 2.87 \times 10^{-4}$, $b = 0.12$).

% === EXAMPLE: TikZ line plot ===
Figure~\ref{fig:scalability-plot} shows the scalability curve.

\begin{figure}[H]
    \centering
    \begin{tikzpicture}[>=Stealth]
        % Axes
        \draw[->] (0,0) -- (8.5,0) node[right] {Records ($\times 10^5$)};
        \draw[->] (0,0) -- (0,5.5) node[above] {Time (s)};

        % Tick marks
        \foreach \x/\val in {1/1, 2/2, 3/3, 4/4, 5/5, 6/6, 7/7, 8/8} {
            \draw (\x,-0.1) -- (\x,0.1) node[below=2pt] {\small \val};
        }
        \foreach \y/\val in {1/50, 2/100, 3/150, 4/200, 5/250} {
            \draw (-0.1,\y) -- (0.1,\y) node[left=2pt] {\small \val};
        }

        % Data points and line (proposed approach)
        \draw[thick, green!60!black, mark=*] plot coordinates {
            (0.1,0.006) (1,0.44) (2,0.88) (4,1.77) (6,2.65) (8,3.53)
        };
        \node[green!60!black, right] at (8,3.53) {\small Proposed};

        % Baseline B
        \draw[thick, blue!60, dashed, mark=square*] plot coordinates {
            (0.1,0.01) (1,0.77) (2,1.55) (4,3.1) (6,4.65) (8,5.2)
        };
        \node[blue!60, right] at (8,5.2) {\small Baseline B};
    \end{tikzpicture}
    \caption{Scalability: processing time vs.\ dataset size} \label{fig:scalability-plot}
\end{figure}

\clearpage
